\uuid{BAto}
\titre{Moment d'ordre 2}
\auteur{Maxime Nguyen}
\organisation{AMSCC}
\datecreate{2026-11-24}
\niveau{L2}
\module{Probabilité}
\chapitre{Probabilité continue}
\sousChapitre{Espérance et Fonction de répartition}
\difficulte{1}
\theme{variance, moment d'ordre 2, formule de König-Huygens}

\contenu{
\begin{question}
\qtype{onlyone}
On note $\mathbb{V}(X)$ la variance d'une variable $X$. La variable $X$ vérifie $\mathbb{E}(X) = 1$ et $\mathbb{V}(X) = 4$. Alors $\mathbb{E}(X^2)$ vaut :
\begin{answers}
\bad{$1$}
\good{$5$}
\bad{$3$}
\bad{$17$}
\end{answers}
\begin{explanations}
Par la formule de König-Huygens : $\mathbb{E}(X^2) = \mathbb{V}(X) + \mathbb{E}(X)^2 = 4 + 1$.

Contre-sens visés : confondre $\mathbb{E}(X^2)$ et $\mathbb{E}(X)^2$ (A) ;
se tromper de sens dans la formule (C) ; confondre variance et écart-type,
c'est-à-dire prendre $\sigma = 4$ et calculer $\sigma^2 + 1 = 17$ (D).
\end{explanations}
\end{question}
}
